% 原文 pp.223–232（印刷 pp.1341–1350）
\section{Lecture 16：四维规范理论}

\begin{keytakeaways}
本讲以欧氏号差的四维规范场论为对象，先从无超对称的纯 Yang--Mills 理论与带 $n_f$ 味的 $SU(N)$ 理论入手，说明红外行为（质量间隙、禁闭）目前仍只能「相信」与「猜测」；核心数学工具是味对称性 $U(n_f)\times U(n_f)$ 的反常破缺、Goldstone 定理，以及流的三点函数所携带的不变三次型 $d^{abc}$——后者给出一个纯粹由短距离计算得到的无质量判据。随后转入 $\mathcal{N}=1$ 超对称纯规范理论，借助超空间表述、全纯性与瞬子，论证手征对称性 $\mathbb{Z}/2h\mathbb{Z}$ 自发破缺出 $h$ 个真空，并由 $\langle\operatorname{Tr}\lambda\lambda\rangle$ 的标度行为反推出 $\beta$ 函数的一圈精确性。最后在 $\mathcal{N}=1$ SQCD 中按 $k=n_f$ 分区讨论红外相，用「下降」程序、量纲与全纯性定出超势 $W=c_{n,k}(\det M)^{-1/(n-k)}$，并指出 $c_{n,n-1}\ne0$ 来自单瞬子计算。全讲的主线是：反常与瞬子这类「短距离/拓扑」信息，足以严格约束甚至决定红外物理。
\end{keytakeaways}

% ===== 原文 p.223 =====

本讲我们将纵览四维规范场论中一些主要的、已知的、被相信的与被猜测的结果，无论是否带有超对称。

\subsection{无超对称的规范理论}
% 原文 §16.1：Gauge theory without supersymmetry

我们从没有超对称（$\mathcal{N}=0$）的情形开始。先考虑单 Lie 群 $G$ 的纯非阿贝尔规范场论（pure nonabelian gauge theory），其拉格朗日量为
\[
L=\frac{1}{4g^2}\int\operatorname{Tr}F\wedge *F .
\]
我们还可以加上带 $\theta$ 角（theta angle）的拓扑项 $\frac{i\vartheta}{16\pi^2}\int\operatorname{Tr}F\wedge F$。这是一个定义良好、渐近自由（asymptotically free）的理论。我们所关心的问题是：这个理论在红外（infrared）端如何行为。

人们相信（对一切 $\vartheta$）这个理论具有质量间隙（mass gap），并猜测它表现出禁闭（confinement）（Wilson 圈算符的面积律，见第 6 讲）。然而，要证明第一个断言、尤其是第二个断言，至今仍是一个重要的未解决问题。

\textbf{Remark.} 正如我们将看到的，超对称的好处在于：在超对称规范场论中，这两个断言至少可以部分地得到证明。

现在让我们加入物质（matter）。本讲采用欧氏号差（Euclidean signature）。先看群为 $G=SU(N)$ 的例子。取物质场为取值于矢量表示 $V$ 与对偶表示 $V^*$ 的手征旋量（chiral spinor）$\psi_i,\psi_i^*$，以及取值于 $V^*$ 与 $V$ 的相应反手征旋量（antichiral spinor）$\bar\psi_i,\bar\psi_i^*$。（手征旋量与反手征旋量分别是取值于给定表示的正、负旋量丛 $S_+$ 或 $S_-$ 的截面。）设每种旋量的个数为 $n_f,n_f^*$（即味（flavor）的数目）。为避免局域反常（local anomaly），必须有 $n_f=n_f^*$（见下一讲）。

带旋量的理论（裸质量为零）的拉格朗日量为
\begin{equation}
  \tag{16.1}
  \label{eq:16-1}
  L=\frac{1}{4g^2}\int\operatorname{Tr}F\wedge *F+\frac{i\vartheta}{16\pi^2}\int\operatorname{Tr}F\wedge F+\sum_i\int\Bigl[(\bar\psi_i,D\psi_i)+(\bar\psi_i^*,D\psi_i^*)\Bigr].
\end{equation}

% [Tutor Note: 记号 $\psi_i^*$ 极易误读。它是与 $\psi_i$ 相互独立的第二个手征旋量（取值于对偶表示 $V^*$），并不是 $\psi_i$ 的厄米共轭；$\psi_i$ 的反手征伙伴是 $\bar\psi_i$。正因如此，$V$ 与 $V^*$ 中的手征旋量味数必须相等（$n_f=n_f^*$），局域反常才能抵消——原文此处说"见下一讲"，即第 17 讲对局域反常的系统讨论。另：本页原文 "The take the matter fields to consist of..." 中 "The take" 显然是 "We take" 的笔误，译文按后者处理。]

% ===== 原文 p.224 =====

\textbf{Remark.} 当 $G=SU(3)$ 且 $n_f=2$ 或 $3$ 时，这就是强相互作用的理论。在这些情形下，若要精确地描述强相互作用，我们其实应当加入小的费米子裸质量。

若 $n_f$ 不太大，该理论是渐近自由的。因此人们关心它的红外行为。

该拉格朗日量具有作用在味上的整体 $U(n_f)\times U(n_f)$ 对称性。然而在量子层面，这个理论有手征反常（chiral anomaly）（因为手征旋量上的 Dirac 算符可以有非零指标（index））。因此，量子力学上算符代数的对称性被破缺到子群 $H=\{(A,B)\in U(n_f)\times U(n_f):\det(A)=\det(B)\}$（群的其余部分被破缺，是因为在瞬子（instanton）场中，它不能使路径积分的测度保持不变）。

人们相信，当 $n_f>1$ 时，群 $H$ 自发破缺（spontaneous symmetry breaking）到 $U(n_f)\times U(n_f)$ 的对角子群 $H_{\mathrm{diag}}$，并由 Goldstone 定理导致 Goldstone 玻色子（Goldstone boson）的存在。注意对角子群正是允许费米子有裸质量的最大子群。若 $n_f\le1$，$H$ 本身已经是对角的，我们就得不到任何自发对称性破缺或 Goldstone 玻色子。在这种情形下，人们相信理论有质量间隙。对 $n_f>1$，Goldstone 玻色子的存在意味着不可能有质量间隙。此时人们相信理论是红外自由的（infrared free），并且除 Goldstone 玻色子之外没有别的无质量粒子；于是低能有效理论是真空空间 $H/H_{\mathrm{diag}}$ 上的一个 $\sigma$ 模型。若再为费米子加上小的裸质量、以使模型更好地描述自然，那么（正如某道习题所探讨的）一般会挑出一个唯一的真空，并由此产生质量间隙。

\textbf{Remark.} 从 $H$ 到对角子群的整体对称性破缺，比我们在第 1 讲中所考虑的情形微妙得多：在目前的情形里，这种破缺在经典层面是看不到的，因为经典真空空间是一个点。就这一点而言，它更类似于二维 $(\bar\psi\psi)^2$ 模型中的手征对称性破缺——后者去年秋天由 Gross 用 $1/N$ 展开分析过。

这个理论的红外行为尚未被完全理解，但某些侧面已经清楚。例如，人们有时可以用下述方式证明质量间隙不存在。下面我们把讨论一般化，考虑任意的四维规范场论，规范群任意紧致、费米子表示任意。令 $H$ 为该理论的（无反常的）整体对称性群。$H$ 的流具有以下有用的性质。

\textbf{Remark.} 下面的分析与我们在第 14 讲开头对二维情形所做的一个分析类似。无论是这里还是第 14 讲，目标都是用短距离计算来了解红外。在第 14 讲中，我们由短距离计算推断了 $\widehat{c}$ 的值；注意 $\widehat{c}$ 非零即意味着不存在质量间隙。这里我们干脆用反常的短距离计算来证明不存在质量间隙。请注意，在二维我们用的是流的二点函数，而在四维我们将用三点函数；在 $2n$ 维的类似论证则会用到 $(n+1)$ 点函数。

% ===== 原文 p.225 =====

令 $J^a$，$a=1,\dots,\dim(H)$，为 $H$ 的流。考虑三点函数 $\langle J^a(x)J^b(y)J^c(z)\rangle$。由于这些流在经典层面守恒，我们得到
\begin{equation}
  \tag{16.2}
  \label{eq:16-2}
  d_x\langle J^a(x)J^b(y)J^c(z)\rangle=0
\end{equation}
只要 $x,y,z$ 互不相同。然而当存在反常时，(16.2) 的左端是支撑在集合 $x=y=z$ 上的非零奇异分布。于是让我们转到动量空间，把 (16.2) 计算到一圈（1-loop）阶。令
\[
F^{abc}_{\lambda\mu\nu}(k,q,r)\,\delta(k+q+r)=\langle \widehat{J}^{\lambda}_a(k)\widehat{J}^{\mu}_b(q)\widehat{J}^{\nu}_c(r)\rangle
\]
为 $J$ 在动量空间的三点函数。于是我们得到（在一圈精度内）
\begin{equation}
  \tag{16.3}
  \label{eq:16-3}
  k_\lambda F^{abc}_{\lambda\mu\nu}=\varepsilon_{\mu\nu\alpha\beta}q^\alpha r^\beta d^{abc},
\end{equation}
其中 $k+q+r=0$，而 $d^{abc}$ 是规范群 Lie 代数上一个不变三次型（invariant cubic form）的分量（$\varepsilon$ 是反对称化张量）。事实上
\[
d^{abc}=\tfrac12\operatorname{Tr}T^a\{T^b,T^c\},
\]
其中 $T^a$ 是与 $J^a$ 相联系的对称性的生成元，作用在手征费米子上。

% [Tutor Note: (16.3) 的指标结构是本节最容易写错的地方：$\lambda$ 是被 $k_\lambda$ 缩并掉的哑指标，$\mu,\nu$ 才是自由指标；右端 $\varepsilon_{\mu\nu\alpha\beta}q^\alpha r^\beta$ 关于 $\mu\leftrightarrow\nu$ 反对称，这正是 $F^{abc}_{\lambda\mu\nu}$ 的反对称性所要求的（若把 $q,r$ 记反，或把 $\varepsilon$ 的指标写在下位，符号就整体出错）。另外，$d^{abc}$ 只取 $\operatorname{Tr}T^a\{T^b,T^c\}$ 的完全对称部分，随后的 $\frac12$ 约定直接影响 (16.10) 中 $c_{n,k}$ 的归一化——跨文献比较系数时必须先核对这个 $\frac12$。]

可以证明这个公式是精确的，没有更高阶修正。事实上，我们在前几讲中已经遇到过张量 $d^{abc}$（例如第 3 讲，以及第 14.1 讲中的二维类比），当时它以「规范化的障碍」的面目出现。$H$ 的那些能被规范化的子群，恰恰是使 $d$ 张量限制到 $H$ 的 Lie 代数上时为零的子群。我们当前公式的一圈精确性，等价于规范化障碍的一圈精确性。

% [Tutor Note: 作者说"It can be proved that this formula is exact"，其背后的定理是 Adler--Bardeen 关于三角图反常系数不受高阶修正的结论：反常是流的三点函数在 $x=y=z$ 处的接触项，其系数由对称性与指标定理锁定（本例中即由 $\operatorname{Tr}T^a\{T^b,T^c\}$ 锁定），故一圈结果精确。作者随后把它与"规范化障碍（obstruction to gauging）的一圈精确性"等同起来，这一步的依据是：若某个子群可被规范化而 $d$ 在其上非零，就会出现 $d$-张量型的局域反常，使路径积分失去规范不变性；而手征流的三点函数在接触项上的系数一旦由对称性固定，高阶修正无法再改变它。]

% [Tutor Note: 此处原文作 "restricted to the Lie algebra of $H$"，按上下文应为"限制到该子群的 Lie 代数"：本段的前提恰恰是 $d$ 在 $H$ 上非零（否则不会有无质量结论），若照字面读成 $d|_{\mathrm{Lie}(H)}=0$，整段论证自相矛盾。]

针对我们当前的目的，关键点如下。任何在 $k=q=r=0$ 附近实解析（real-analytic）的函数 $F$ 都不可能以满足 (16.3) 的方式（且 $d$ 非零）存在。因此若 $d$ 不为零，$F$ 在零动量附近是奇异的。这意味着：若 $d\ne0$，理论中存在无质量粒子，因而没有质量间隙。

于是，在某些情形下，我们可以仅凭一个短距离计算就看出四维规范场论没有质量间隙。

我们希望更清楚地理解 $d^{abc}\ne0$ 诸情形的红外极限。我们回忆第 2 讲第 4、5 节：四维的红外自由临界点受到严重限制。无质量粒子的自旋只能是 $0$、$1/2$ 或 $1$。在没有超对称时，人们有理由预期：自旋 $0$ 的无质量粒子是 Goldstone 玻色子，而自旋 $1/2$ 的无质量粒子之所以无质量，是因为存在未破缺的手征对称性。

让我们看看 $d^{abc}\ne0$ 时有哪些可能性。

\textbf{可能性 1.} 理论在红外不是自由的。除下文关于它在何时出现的评论之外，我们不追究这种可能性。

\textbf{可能性 2.} 理论在红外是自由的。于是 Green 函数可以用轻场的 Feynman 图计算。因此必然存在对 $F$ 的非解析性有贡献的图。这又有两种可能的情形：

% ===== 原文 p.226 =====

\textbf{情形 1.} 没有 Goldstone 玻色子：场 $J(x)$ 不能从真空中产生单粒子的无质量态（$J$ 的二点函数在零动量处没有极点）。那么张量 $d^{abc}$ 只能来自圈图。这类图内部必须有质量为零的费米子，因为玻色子的一圈图不产生反常。这些无质量费米子未必是原拉格朗日量中的那些——例如它们未必按同样的方式在 $H$ 下变换——但它们对 $d$ 的贡献必须与短距离拉格朗日量中的费米子相同。这一可能性是 't Hooft 在七十年代后期指出的，但直到最近几年才借助超对称找到例子。欲知更多，请参阅 Seiberg 的讲义。

\textbf{情形 2.} 没有无质量费米子，因而没有对 $d$ 的圈图贡献。此时 $d$ 张量必须来自树图。这只有在 $J$ 的二点函数于零动量处有极点时才可能，对应于 $J$ 从真空中产生一个 Goldstone 玻色子这一事实（见第 1 讲）。该态是自旋为 $0$ 的无质量玻色子 $\pi$，使得在红外极限下 $J=*d\pi$。这个玻色子 $\pi$ 就是被破缺的整体 $H$ 对称性的 Goldstone 玻色子。若存在由 $J^a$ 产生的 Goldstone 玻色子，并且还存在合适的非零矩阵元 $\langle\pi|J^bJ^c|\Omega\rangle$（使两个流的乘积能从真空态 $|\Omega\rangle$ 产生 $|\pi\rangle$），那么就可以构造一个带 Goldstone 玻色子交换的树图，由它生成 $d$ 张量。

\textbf{Remark.} 以上表述有些过于简化。例如，可以有情形 1 与情形 2 的混合：$H$ 自发破缺到子群 $H'$，使得 $d$ 限制到 $H'$ 的 Lie 代数上时非零，同时存在处于 $H'$ 的手征表示中的无质量费米子。

\textbf{Remark.} 即使在情形 1 中，二点函数也不必在原点解析。若没有质量间隙，则 $J$ 的二点函数在零动量处几乎总是不解析的：它有一条从零动量开始的支割（branch cut）。

于是结论是：若理论没有无质量费米子，则 $H$ 对称性自发破缺到某个群 $H'\subset H$，使得 $\left.d\right|_{H'}=0$。

这正是自然界中所发生的事。也就是说，真空只在对角子群下不变；对该群而言 $d=0$，因为费米子所在的表示在这个群下是实的，从而所有手征场都被反手征场配平。这是 $H$ 中最显然的满足 $d=0$ 的子群。

\textbf{Remark.} 就所涉及的物理而言，情形 2 的确相当富有戏剧性。所谈的树图描述的是中性 $\pi$ 介子衰变为两个光子。正是在研究这一衰变时，S. Adler、J. Bell 与 R. Jackiw 于 1970 年前后发现的反常。

现在我要谈谈可能性 1——非自由的红外极限——在什么条件下是可预期的。为说明起见，回到我们最初考虑的例子——$SU(n)$ 规范理论，费米子由 $n_f$ 份 $V\oplus V^*$ 组成。回忆渐近自由的条件是 $n_f<11n/3$。若 $n_f$ 远小于 $11n/3$，人们相信理论是红外自由的，行为如上所述（$H$ 自发破缺到无反常的对角子群；除 Goldstone 玻色子外没有无质量粒子）。人们还预期，当 $n_f/n$ 足够小时理论表现出禁闭。然而，若 $n_f$ 接近 $11n/3$，则预期理论流向一个非平凡的、红外稳定的不动点。这

% ===== 原文 p.227 =====

是因为这些理论 $\beta$ 函数的二圈修正为正，即使在 $n_f=11n/3$ 时也是如此。因此，若取比值 $n_f/(11n/3)$ 略小于 $1$、使一圈 $\beta$ 函数变得很小，则 $\beta$ 函数在原点附近、且在微扰论可及的范围内有一个红外稳定的零点。当然，既然 $n_f$ 是整数，要使 $n_f/(11n/3)$ 略小于一，就必须同时取 $n$ 很大；所以这一现象也许只在 $n$ 足够大时才出现。在相应的超对称情形中，人们确切知道 $n$ 与 $n_f/n$ 要多大才能得到非自由的红外不动点；这在 Seiberg 的讲义中有解释。

\subsection{$\mathcal{N}=1$ 超对称纯规范理论}
% 原文 §16.2：$\mathcal{N}=1$ supersymmetric pure gauge theory

现在考虑 $\mathcal{N}=1$ 的超对称理论。我们从「纯规范理论」开始。用分量场来说，该理论有一个规范场 $A$ 和一个取值于伴随表示的手征旋量 $\lambda$。拉格朗日量为
\[
L=\frac{1}{4g^2}\int\operatorname{Tr}F\wedge *F+\frac{i\vartheta}{16\pi^2}\int\operatorname{Tr}F\wedge F+\int\operatorname{Tr}\bar\lambda\, iD\lambda .
\]

这个理论在经典层面的对称性是手征 $U(1)$ 对称性 $\lambda\to e^{i\delta}\lambda$。然而，由于伴随丛中旋量上的 Dirac 算符 $D$ 可以有非平凡指标，这个手征对称性在量子层面并不存在。不过，可以证明这个指标在球面 $S^4$ 上总是被 $2h$ 整除，其中 $h$ 是该 Lie 代数的对偶 Coxeter 数（对第二 Chern 类为 $k$ 的丛，它等于 $2hk$）。这表明（若在 $\mathbb{R}^4$ 或 $S^4$ 上）$U(1)$ 中由积分测度的对称性所组成的子群是 $\mathbb{Z}/2h\mathbb{Z}$。这个群是该理论的对称群；它作用在量子算符代数上。

在红外极限下，群 $\mathbb{Z}/2h\mathbb{Z}$ 有可能自发破缺到某个子群。它能破缺到的最小子群是 $\{-1,1\}$，因为 $-1$ 是 Poincaré 群的双覆盖的中心元素。

关于这个理论的标准预期是：
\begin{enumerate}
  \item 质量间隙。
  \item 禁闭。
  \item 手征对称性破缺：$\mathbb{Z}/2h\mathbb{Z}$ 自发破缺到 $\{1,-1\}$。恰好有该对称性破缺所要求的 $h$ 个真空。
\end{enumerate}

第一个断言意味着费米子质量是动力学地产生的。这暗示 $\mathbb{Z}/2h\mathbb{Z}$ 被破缺到的子群必须允许基本费米子有裸质量。最大的这种子群是 $\mathbb{Z}/2\mathbb{Z}$。所以断言 1 暗示断言 3。

值得注意的是，关于对称性破缺的断言 3 实际上是可以检验的。为此，只要证明算符 $\operatorname{Tr}\lambda\lambda$ 的一点函数 $\langle\operatorname{Tr}\lambda\lambda\rangle$ 不为零即可（所谓 $\operatorname{Tr}\lambda\lambda$，指的是把 $\lambda$ 的正手征分量用旋量上的斜对称形式组合起来、并在 Lie 代数上取迹所得到的 Lorentz 标量算符）。

若恰好有 $n$ 个被破缺的手征对称性所置换的真空，则在第 $k$ 个真空中 $\langle\operatorname{Tr}\lambda\lambda\rangle=Ce^{2\pi ik/h}$，其中 $C$ 为某个常数。因为群 $\mathbb{Z}/2h\mathbb{Z}$ 的生成元把第 $k$ 个真空送到第 $k+1$ 个，并把 $\langle\operatorname{Tr}\lambda\lambda\rangle$ 乘以 $e^{2\pi i/h}$。

% [Tutor Note: 这里的"$n$ 个真空"应为 $h$ 个。$\mathbb{Z}/2h\mathbb{Z}$ 破缺到 $\mathbb{Z}/2\mathbb{Z}$ 恰好给出 $h$ 个真空（本页上文与 p.229 的 "the $h$-point function of $\operatorname{Tr}\lambda\lambda$" 都可印证），而相位 $e^{2\pi ik/h}$ 也以 $h$ 为周期；$n$ 在 $\mathcal{N}=1$ 纯规范理论里根本没有定义（$n$ 是下一节 $SU(n)$ 的秩）。另注意 $k$ 取值 $0,\dots,h-1$，故 $\langle\operatorname{Tr}\lambda\lambda\rangle$ 是 $\vartheta$ 的 $h$ 值多值函数，只有它的 $h$ 次幂才单值。]

% ===== 原文 p.228 =====

现在让我们自问：函数 $\langle\operatorname{Tr}\lambda\lambda\rangle$ 如何依赖于 $\theta$ 角？令 $J$ 为 $U(1)$ 对称性的流。由于手征反常，$dJ$ 非零，且等于 Chern--Weil 形式的一个倍数：
\begin{equation}
  \tag{16.4}
  \label{eq:16-4}
  dJ=\frac{h}{8\pi^2}\operatorname{Tr}F\wedge F .
\end{equation}

这意味着把 $\lambda$ 作 $e^{i\alpha}$ 的转动，等价于给 $\theta$ 角加上 $2h\alpha$：$\vartheta\to\vartheta+2h\alpha$（也就是说，$U(1)$ 在理论空间上的这两个作用是重合的）。因此我们应预期
\[
\langle\operatorname{Tr}\lambda\lambda\rangle=Ce^{i\vartheta/h},
\]
其中 $C$ 为实常数且与 $\theta$ 角无关。注意这个函数是多值的，有 $h$ 个不同的值，对应于 $h$ 个真空。

既然 $\langle\operatorname{Tr}\lambda\lambda\rangle$ 有 $h$ 个不同的值，取它的 $h$ 次幂 $\langle\operatorname{Tr}\lambda\lambda\rangle^h=C^he^{i\vartheta}$ 较为方便。这是 $\vartheta$ 的单值函数。

常数 $C$ 只依赖于规范耦合 $g$，且在微扰论中为零，因为微扰论在整个 $U(1)$ $R$ 对称性下不变，而不只是在无反常的 $\mathbb{Z}/2h\mathbb{Z}$ 下不变。因此，若 $C\ne0$，它是一种非微扰（nonperturbative）现象（就像 Gross--Neveu 模型中那样，见 Gross 的第 4 讲）。

到此为止我们所做的，只是说明：若 $\mathbb{Z}/2h\mathbb{Z}$ 手征对称性自发破缺，则必然发生什么。现在让我们计算 $\langle\operatorname{Tr}\lambda\lambda\rangle^h$。我们有
\begin{equation}
  \tag{16.5}
  \label{eq:16-5}
  \langle\operatorname{Tr}\lambda\lambda\rangle^h=\lim_{|x_i-x_j|\to\infty}\langle\operatorname{Tr}\lambda\lambda(x_1)\dots\operatorname{Tr}\lambda\lambda(x_n)\rangle
\end{equation}
（团簇分解（cluster decomposition））。这个断言在所有的真空中都成立。

% [Tutor Note: (16.5) 左端是 $h$ 次幂，右端却写成 $n$ 个算符的乘积，此处的 $n$ 应为 $h$。团簇分解给出的是 $\langle(\operatorname{Tr}\lambda\lambda)^h\rangle=\lim_{|x_i-x_j|\to\infty}\langle\operatorname{Tr}\lambda\lambda(x_1)\cdots\operatorname{Tr}\lambda\lambda(x_h)\rangle$；若照字面取 $n$，则 $n$ 在本节中从未定义，且 p.229 的推理链（"$h$ 点函数非零 $\Rightarrow\langle\operatorname{Tr}\lambda\lambda\rangle\ne0$"）无法闭合。]

我们想证明 $\langle\operatorname{Tr}\lambda\lambda(x_1)\dots\operatorname{Tr}\lambda\lambda(x_n)\rangle$ 与 $x_1,\dots,x_n$ 无关。为此，回忆本模型的超空间（superspace）表述（见 Superhomework 与 Supersolutions）。

在超空间表述中，我们在超时空（super-spacetime）$\mathbb{R}^{4,4}$ 上描述这个理论，其坐标为 $x_0,\dots,x_3$ 以及 $\theta_+,\theta_-,\bar\theta_+,\bar\theta_-$。本模型的拉格朗日量等价于拉格朗日量
\[
L'=\int d^4x\bigl(d^2\theta\,\tau\operatorname{Tr}WW+d^2\bar\theta\,\bar\tau\operatorname{Tr}\bar W\bar W\bigr),
\]
其中 $\tau=\frac{\vartheta}{2\pi}+\frac{4\pi i}{g^2}$，而 $W$ 是超曲率（supercurvature）：$W=\lambda+\Theta F_++\dots$，$\bar W=\bar\lambda+\Theta F_-+\dots$。这里 $\Theta$ 是 $\mathbb{R}^{4,4}$ 上取值为 $S_+$ 的规范线性手征函数，$F_+$ 是曲率的自对偶部分，而 $\operatorname{Tr}W^2$ 是用旋量上的斜对称形式缩并旋量 $W$ 所得到的标量。$W$ 的公式中 $\Theta F_+$ 这一项涉及 $SU(2)$ 映射 $\mathbb{C}^2\otimes\mathbb{C}^3\to\mathbb{C}^2$。

超曲率是取值于旋量的手征超场。具体说，令 $\bar D$ 为湮灭手征函数的、取值于旋量的算符（我们有 $[\bar D,\bar D]=0$）。那么 $\bar DW=0$。这意味着 $\bar D(\operatorname{Tr}WW)=0$。由于 $d=[D,\bar D]$，我们得到 $d(\operatorname{Tr}W^2)=\bar D\bar D(\operatorname{Tr}W^2)$。

现在回忆：对任何算符 $X$，一点函数 $\langle\bar DX\rangle$ 为零。（这个断言反映了路径积分测度的 $\bar D$ 不变性，但愿大家从第 14、15 讲已经熟悉。）利用 $\bar D(\operatorname{Tr}W^2)=0$ 与 $d(\operatorname{Tr}W^2)=\bar D(\dots)$，这意味着 $\langle d(\operatorname{Tr}WW(x_1))\operatorname{Tr}WW(x_1)\dots\operatorname{Tr}WW(x_n)\rangle=0$。在取 $\Theta=0$ 之后，这意味着 $\langle\operatorname{Tr}\lambda\lambda(x_1)\dots\operatorname{Tr}\lambda\lambda(x_n)\rangle$ 与 $x_1,\dots,x_n$ 无关。

% [Tutor Note: $d=[D,\bar D]$ 是简写，省略了 $\sigma$ 矩阵的缩并。在 $\mathcal{N}=1$ 超空间中 $\{D_\alpha,\bar D_{\dot\beta}\}=-2i\sigma^\mu_{\alpha\dot\beta}\partial_\mu$，须再乘以 $\frac{i}{2}(\sigma_\mu)^{\alpha\dot\beta}$ 并对自旋指标缩并，才得到 $\partial_\mu$ 进而得到 $d$；同一个式子的 $[\bar D,\bar D]=0$ 说的是同手征分量的分次对易关系 $\{\bar D_{\dot\alpha},\bar D_{\dot\beta}\}=0$。另外，$\Theta F_+$ 这一项之所以是 $\mathbb{C}^2\otimes\mathbb{C}^3\to\mathbb{C}^2$ 的 $SU(2)$ 不变映射，是因为自对偶二形式空间即 $\mathrm{Sym}^2S_+\cong\mathbb{C}^3$；把 $F_+$ 的指标与 $\Theta$ 的自旋指标配错，就会得到错误的 $F_-$ 分量。]

% ===== 原文 p.229 =====

这意味着我们可以在短距离处计算 $\langle(\operatorname{Tr}\lambda\lambda)^h\rangle$，也就是假定所有 $x_i$ 之间的间距都远小于重整化理论时所引入的长度标度（这个长度标度是 $1/\mu$，其中 $\mu$ 是重整化中引入的质量标度）。在短距离处，$\langle(\operatorname{Tr}\lambda\lambda)^h\rangle$ 可以用渐近自由来计算。另一方面，由 (16.5)，既然 $\operatorname{Tr}\lambda\lambda$ 的 $h$ 点函数非零，它在短距离处取非零值就意味着 $\langle\operatorname{Tr}\lambda\lambda\rangle$ 非零。一旦有一个真空具有这一性质，就必然至少有 $h$ 个真空被 $\mathbb{Z}/2h\mathbb{Z}$ 所置换。这就是我们想要建立的结构。

我们的参数是 $\mu$（重整化标度）、$g$（规范耦合）与 $\vartheta$（$\theta$ 角）。于是 $\langle\operatorname{Tr}\lambda\lambda\rangle=F(\mu,g,\vartheta)$。由量纲论证，
\[
F(\mu,g,\vartheta)=\mu^{3h}f(g,\vartheta),
\]
其中 $f$ 无量纲（四维中费米子的标度维数为 $3/2$）。此外我们已发现 $\langle\operatorname{Tr}\lambda\lambda\rangle^h$ 具有形式 $Ce^{i\vartheta}$，其中 $C$ 与 $\vartheta$ 无关。这意味着 $f(g,\vartheta)=\rho(g)e^{i\vartheta}$。

现在到了要点。我们的理论有一个由超对称导出的关键性质：若算符 $X$ 满足 $\bar DX=0$，则 $\langle X\rangle$ 是 $\tau$ 的全纯函数，其中 $\tau$ 是超空间表述中的「模」参数。这来自如下事实：
\[
\frac{\partial L'}{\partial\bar\tau}=\int d^4x\,d^2\bar\theta\operatorname{Tr}\bar W^2=\bar DY,
\]
其中 $Y=\bar D\operatorname{Tr}\bar W^2$。但对这样的 $Y$ 与手征的 $X$，有 $\langle YX\rangle=\langle\bar D(\bar D\operatorname{Tr}\bar W^2\cdot X)\rangle=0$。因此 $F$ 关于 $\tau$ 全纯。这意味着 $f(g)=e^{-8\pi^2/g^2}$，即
\begin{equation}
  \tag{16.6}
  \label{eq:16-6}
  F(\mu,g,\vartheta)=C_0\exp\Bigl(i\vartheta-\frac{8\pi^2}{g^2}\Bigr),
\end{equation}
其中 $C_0$ 是常数。

% [Tutor Note: 由 (16.6) 与 (16.7) 可以干净地反解出 $\beta$。(16.6) 中隐含因子 $\mu^{3h}$（$\langle\operatorname{Tr}\lambda\lambda\rangle^h$ 的标度维数为 $3h$，原文未显式写出、被吸收进 $C_0$），故 $\mu\partial_\mu F=3hF$，而 $\beta\partial_gF=\beta F\cdot\frac{16\pi^2}{g^3}$；代入 (16.7) 得 $3h+\beta\cdot16\pi^2/g^3=0$，即 $\beta=-3hg^3/(16\pi^2)$。]

% [Tutor Note: 常数 $C_0$ 只能来自单瞬子，理由是零模计数：$\langle\operatorname{Tr}\lambda\lambda\rangle^h$ 在 $R$ 对称性下带荷 $2h$，故所涉及的 Dirac 算符指标必须为 $2h$；而瞬子数为 $1$、取值于伴随表示的旋量恰好有 $2h$ 个零模（指标 $=2hk$，$k=1$），正好把 $2h$ 个 $\lambda$ 插入饱和掉。$k\ge2$ 时零模数为 $2hk>2h$，无法被 $\langle\operatorname{Tr}\lambda\lambda\rangle^h$ 饱和，故不贡献。]

常数 $C_0$ 不易求得。要求 $C_0$，应注意 $\langle\operatorname{Tr}\lambda\lambda\rangle^h$ 的唯一非平凡贡献只能来自瞬子数（即第二 Chern 类）等于 $1$ 的场构型。$C_0$ 只能来自瞬子数 $1$ 的理由是：$\langle\operatorname{Tr}\lambda\lambda\rangle^h$ 在 $R$ 对称性下带荷 $2h$，所以 Dirac 算符的指标必须是 $2h$。另一个有吸引力的检验是：$8\pi^2/g^2$ 正是瞬子作用量，这解释了为何 $g$ 依赖是 $\exp(-8\pi^2/g^2)$。于是，常数 $C_0$ 可以通过在单个瞬子附近做微扰论来求得。特别地，可以用一个相当复杂的计算证明 $C_0\ne0$（动力学对称性破缺）。

\textbf{Remark 1.} 由于 $\bar W=\bar\lambda+\Theta F_-+\dots$，且瞬子由「曲率自对偶且 $\lambda=0$」定义，故对瞬子有 $\bar W=0$，从而拉格朗日量在瞬子处（且仅在瞬子处）关于 $\tau$ 全纯。这提示了为什么函数 $F$ 应当从瞬子获得贡献。

\textbf{Remark 2.} 算符 $F$ 的反常维数为零，因为例如它的一个分量就是积分后给出规范丛第二 Chern 类的密度；后者是整数，不受量子重整化影响。作为规范不变算符的期望值，$F$ 满足重整化群方程
\begin{equation}
  \tag{16.7}
  \label{eq:16-7}
  \Bigl(\mu\frac{\partial}{\partial\mu}+\beta(g)\frac{\partial}{\partial g}\Bigr)F=0,
\end{equation}
其中 $\beta$ 是 $\beta$ 函数。先验地我们只知道 $\beta(g)=-3hg^3+O(g^5)$。公式 (16.6) 表明一圈 $\beta$ 函数其实是精确的：
\begin{equation}
  \tag{16.8}
  \label{eq:16-8}
  \beta(g)=-3hg^3 .
\end{equation}

% [Tutor Note: (16.8) 与 p.230 的 $\beta=-\frac{g^3}{\pi}(3n-k)$ 都未写 $16\pi^2$。若采用使瞬子作用量为 $8\pi^2/g^2$（即 (16.6) 的指数）的标准归一化，由 (16.6)、(16.7) 反解得到的是 $\beta=-3hg^3/(16\pi^2)$；这两处要么整体省略了 $16\pi^2$（等价于把耦合相对标准约定重新标度为 $g_{\mathrm{std}}=4\sqrt{\pi}\,g$），要么与 (16.6) 的归一化不自洽。更麻烦的是这两页彼此也差一个 $\pi$：令 $k=0$、$h=n$，p.229 的 (16.8) 给 $-3ng^3$，而 p.230 给 $-3ng^3/\pi$。引用这些系数前务必先统一约定。]

% ===== 原文 p.230 =====

这等于说，参数 $\tau$ 只有加性重整化。在微扰论中，无论如何这一点由全纯性就很明显：用 $g^2$ 的非恒定幂级数去重整化 $\tau$，必定破坏全纯性。$\tau$ 的加性重整化取形式
\[
\tau(\mu)=\tau(\mu_0)+48\pi^2h\ln(\mu/\mu_0).
\]

% [Tutor Note: 这个加性系数几乎肯定是笔误。由 $\tau=\frac{\vartheta}{2\pi}+\frac{4\pi i}{g^2}$ 与 $\beta=-3hg^3$ 可得 $\frac{d}{d\ln\mu}\bigl(\frac{1}{g^2}\bigr)=6h$，故 $\tau(\mu)=\tau(\mu_0)+24\pi i\,h\ln(\mu/\mu_0)$（若改用带 $16\pi^2$ 的 $\beta$，系数为 $\frac{3h}{2\pi}i$）。两种算法给出的系数都是纯虚数，这是因为 $\operatorname{Re}\tau=\vartheta/2\pi$ 在微扰论中不跑动；而印出的实数 $48\pi^2h$ 只会平移 $\vartheta$ 而不改变耦合，与"加性重整化 $\tau$"的说法直接冲突。疑为脱漏因子 $i$、并把 $\pi$ 误排成 $\pi^2$。]

\subsection{带手征超场的 $\mathcal{N}=1$ 理论}
% 原文 §16.3：$\mathcal{N}=1$ theories with chiral superfields

现在让我们考虑同一个理论再加上一个取值于伴随表示的手征超场 $\Phi$。新的拉格朗日量为
\begin{equation}
  \tag{16.9}
  \label{eq:16-9}
  L''=L'+\int d^4x\,d^4\theta\operatorname{Tr}\bar\Phi\Phi+\Bigl(\varepsilon\int d^4x\,d^2\theta\operatorname{Tr}\Phi^2+\text{c.c.}\Bigr).
\end{equation}
（c.c. 即「复共轭」。）这里 $\varepsilon$ 是一个复数，充当 $\Phi$ 的质量参数。若 $\varepsilon$ 非常大，把 $\Phi$ 积掉之后，我们就重新得到先前的理论作为有效理论。若 $\varepsilon=0$，我们得到 $\mathcal{N}=2$ 超 Yang--Mills 理论，它与 Donaldson 理论有密切关系。

\textbf{Remark.} 由超对称（在标准假设之下）可以推知：任何非零的 $\varepsilon$，无论多小，其情形都与大 $\varepsilon$ 的情形在红外的定性行为上相同。证明的办法是把 $\Phi$ 重新定义为 $\Phi\to\Phi/\sqrt{\varepsilon}$。这一操作把 $\varepsilon$ 从超势中移走。此时 $\varepsilon$ 的依赖只出现在 $\Phi$ 的动能项里，而那是一项形如 $\int d^4\theta(\dots)$ 的项。人们相信，这类项在超对称理论的红外行为中是普遍无关的（irrelevant），因此预期本理论的定性行为与 $\varepsilon$ 无关。

上述 Remark 意味着：对 $G=SU(2)$，当 $\varepsilon\ne0$ 时本理论有两个真空（与没有超场 $\Phi$ 的理论相同）。这只要取大 $\varepsilon$ 即可看出：此时 $\Phi$ 极重，因而无关紧要。另外，我们将在第 18 讲证明：当 $\varepsilon$ 小而非零时，本理论有质量间隙与禁闭。于是由上面的 Remark，这些性质对 $\varepsilon\to\infty$ 成立，从而对纯 $\mathcal{N}=1$ 超对称规范理论也成立。

对 $\varepsilon=0$，行为完全不同。存在一族单复参数的真空，由 $u=\operatorname{Tr}\Phi^2$ 参数化，我们下周将讨论。

现在我们要考虑一个很不相同的例子。取 $G=SU(n)$，讨论 $\mathcal{N}=1$ 超对称 QCD。这意味着我们去掉取值于伴随表示的手征超场，转而考虑处于表示 $R=V^{\oplus n_f}\oplus V^{*\oplus n_f}$ 的手征超场 $\Phi$，其中 $V$ 是 $SU(n)$ 的矢量表示。（这正是本讲开头所用的例子；和以前一样，$V$ 与 $V^*$ 的重数必须相同以避免反常。）我们把超场写成 $\Phi=(Q_1,\dots,Q_k,Q_1^*,\dots,Q_k^*)$，其中 $k=n_f$。该理论渐近自由当且仅当 $k<3n$，因为 $\beta=-\frac{g^3}{\pi}(3n-k)$。

这个理论的动力学行为非常不平凡。更详细的情形见 [II-N=1 Dynamics]。人们相信下面的图像成立：
\begin{enumerate}
  \item $k=0$——上文已讨论。有质量间隙。
  \item $k=1,\dots,n-1$——下文讨论。低能理论是带非零超势的超对称 $\sigma$ 模型。除非为手征多重态加上裸质量，否则根本不存在稳定真空。

% ===== 原文 p.231 =====

  \item $k=n,n+1$——理论是红外自由的，并且有一个有趣的真空，其中的无质量场比经典上所能预期的更多，或至少不同。
  \item $n+2\le k<3n/2$——理论是红外自由的，且在红外由一个规范理论描述，但红外的规范群不是紫外规范群的子群。
  \item $3n/2\le k<3n$——红外极限是非平凡的、非自由的。
\end{enumerate}

今天我们只讨论 $1\le k\le n-1$ 的情形。我们将试图解释超势为何非零。

让我们计算经典的真空模空间（moduli space of vacua）。按第 2 讲，这个模空间应当是辛商（symplectic quotient）：$\mathbf{M}=R//G$（其中 $R$ 和以前一样是手征超场表示）。这里 $R$ 上的辛结构由 $R$ 上的厄米形式诱导。由于 $G$ 全纯地作用，我们有 $R//G=R/G_{\mathbb{C}}$。于是我们化归为经典不变量理论。

对 $k<n$，$R$ 中 $SL(n)$ 的所有不变量都是内积 $M_{ij}=(Q_i,Q_j^*)$ 的函数。这些不变量彼此独立，并在 $R/G_{\mathbb{C}}$ 上构成坐标系。此外，让我们计算 $\det(M)\ne0$ 的点的稳定化子（stabilizer）。容易看出这个稳定化子同构于 $SL(n-k)$。于是我们得到一个 Higgs 现象：$SU(n)$ 规范对称性自发破缺到 $SU(n-k)$（见第 2 讲 §2.2）。特别地，当 $k=n-1$ 时它被完全破缺。这说明对 $k<n-1$，经典理论有 $SU(n-k)$ 的无质量规范场，以及参数化辛商的无质量手征多重态。量子力学上，事情会因为若干原因而变得更复杂。

特别地，对 $n-k\ge2$，由于 $\mathbf{M}$ 中某点的稳定化子是 $SU(n-k)$，我们应预期低能有效理论是一个 $\mathcal{N}=1$ 超对称的 $SU(n-k)$ 规范理论（外加在 $SU(n-k)$ 下中性的轻手征超场）。正如我们在今天这讲第 2 节中所讨论的，人们相信这个理论有 $n-k$ 个真空，因此我们应预期每个经典真空对应 $n-k$ 个量子真空。换句话说，量子真空空间将是 $\mathbf{M}$ 的一个 $(n-k)$ 重覆盖。称这个覆盖为 $\widetilde{\mathbf{M}}$。于是低能有效理论涉及一个映入 $\widetilde{\mathbf{M}}$ 的 $\sigma$ 模型。这样的超对称 $\sigma$ 模型有可能容许一个超势。与上面的论证类似可以证明：这个超势 $W$ 必须是 $\widetilde{\mathbf{M}}$ 上的全纯函数，并且作为复化规范耦合 $\tau$ 的函数也是全纯的。超势当然应当具有与理论本身相同的对称性。

经典地，我们的理论具有作用在超场上的 $U(k)\times U(k)$ 整体对称性，以及 $R$ 对称性——即转动 $\lambda$ 的 $U(1)$ 对称性。这些对称性中有一部分是反常的（手征反常）。作用在量子理论上的无反常子群由群 $H=\{(A,B)\in U(k)\times U(k);\det(A)=\det(B)\}$ 与群 $U(1)=\{(\zeta,\zeta^{-(n-k)/n},\zeta^{(n-k)/n})\in U(1)\times U(k)\times U(k)\}$ 生成。由这些对称性（并回忆本讲第一部分中 $R$ 对称性如何作用在 $\theta$ 角上），我们看到超势必须等于
\begin{equation}
  \tag{16.10}
  \label{eq:16-10}
  W(M)=c_{n,k}(\det M)^{\frac{-1}{n-k}} .
\end{equation}

于是我们确实得到了 $\mathbf{M}$ 的 $(n-k)$ 重覆盖 $\widetilde{\mathbf{M}}$ 上的一个函数，正如所预期。

现在我们需要计算 $c_{n,k}$，特别是要看它是否为零。做到这一点的关键事实是：若 $c_{n,k}\ne0$，则 $c_{n-1,k-1}$ 非零，且其中一个可以由另一个确定。

% ===== 原文 p.232 =====

事实上，若在 $(n,k)$ 理论中取一个具有一个非常大本征值的 $M$，则规范群从 $SU(n)$ 自发破缺到 $SU(n-1)$，一些手征超场获得质量，而我们实际上被约化为把 $n,k$ 换成 $n-1,k-1$ 的理论。

于是若 $k<n-1$，我们可以从 $(n,k)$ 下降（descend）到 $(n-k,0)$，从而把问题归结为纯 $\mathcal{N}=1$ 规范理论的情形。在纯规范理论中，空间 $\mathbf{M}$ 是 $n-k$ 个点的并，而超势 $W$ 的角色由我们已经算过的对象 $\langle\operatorname{Tr}\lambda\lambda\rangle$ 扮演。因此 $c_{n-k,0}\ne0$。

\textbf{Remark.} 看起来超势似乎只能定义到「加上一个局部常函数」为止。然而事实上它只定义到「加上一个常函数」为止。这是因为在所讨论的空间紧致的情形下只有一个真空，于是模空间的 $h$ 个点彼此「知道」对方；可以在它们之间作插值（通过一族并非真空的物理态），从而比较不同真空处超势的值。

现在令 $k=n-1$。于是下降程序将停在 $n=2,k=1$。在 $n=2,k=1$ 的情形下，有 $W=c_{2,1}(\det(M))^{-1}$。更精确地说，由量纲论证 $c_{2,1}=\widetilde{c}_{2,1}\mu^5$，其中 $\mu$ 是该理论的重整化标度，$\widetilde{c}_{2,1}$ 无量纲。正如纯规范理论中对 $\langle\operatorname{Tr}\lambda\lambda\rangle$ 的讨论那样，可以证明对 $c_{2,1}$ 有贡献的瞬子具有瞬子数 $1$。于是由 $W$ 关于 $\tau$ 的全纯性，我们有
\begin{equation}
  \tag{16.11}
  \label{eq:16-11}
  W=\widetilde{c}_{2,1}\mu^5e^{i\vartheta-8\pi^2/g^2}/\det(M).
\end{equation}

$\widetilde{c}_{2,1}$ 的计算很棘手。原因在于存在各种大小的瞬子，这使得在瞬子周围做微扰论更为困难。

这个计算由 Affleck、Dine 与 Seiberg 完成。他们证明了 $c_{2,1}$ 不为零。因此 $c_{n,n-1}$ 不为零。

关于上述内容、以及 $k\ge n$ 时会发生什么，我请读者参阅 Seiberg 的讲义。

\subsection*{参考文献与交叉引用}

关于四维场论红外极限的一般讨论，见第 2 讲。四维 $\mathcal{N}=1$ 规范理论的拉格朗日量在第 8 章 [I-Supersolutions] 中讨论。§16.3 的讨论可看作 Seiberg 讲义 [II-N=1 Dynamics] 的引论；那里详细解释了这些理论，并可找到更多文献。
